% NichidaiMasterExam_R7_1st_4
\begin{tcolorbox} %%R7_1st_4
	$\fbox{4}$ $a, b$を正の定数とする。3次元空間$\R^3$において、$z= x^2+ 2y^2$を$S$として$(x_0, y_0, z_0)$を$S$上の点とする。
\begin{enumerate}
	\item\label{vol:楕円柱} 曲面$S$、平面$z= 0$、楕円柱$\frac{x^2}{a^2}+ \frac{y^2}{b^2}= 1$で囲まれた立体の体積を求めよ。
	\item $a^2 +2b^2= 1$において\ref{vol:楕円柱}の立体の体積の最大値を求めよ。
	\item 曲面$S$の円柱$x^2+ 4y^2= 1$で囲まれた部分の曲面積を求めよ。
	\item 曲面$S$の$(x_0, y_0, z_0)$における接平面の方程式を求めよ。
\end{enumerate}
\end{tcolorbox}

\begin{souanswer} %%R7_1st_4_ans
	$\fbox{4}$
\begin{enumerate}
	\item 楕円柱$\frac{x^2}{a^2}+ \frac{y^2}{b^2}\le 1$を底面$D$とし、高さ$z= x^2+ 2y^2$の立体の体積を求める。
\begin{equation*}
    V= \iint_D (x^2+ 2y^2)dxdy
\end{equation*}
    楕円座標$x= ar\cos{\theta}, y= br\sin{\theta}$を用いて置換積分を行う。ヤコビアンは$J= abr$で積分範囲は$0\le r\le 1, 0\le \theta\le 2\pi$。
\begin{align*}
    V
    &= \int_{\theta= 0}^{2\pi}\int_{r= 0}^1 (a^2r^2\cos^2{\theta}+ 2b^2r^2\sin^2{\theta})\cdot abrdrd\theta\\
    &= ab\int_{\theta= 0}^{2\pi}\left[\frac{r^4}{4}\right]_{r= 0}^1(a^2\cos^2{\theta}+ 2b^2\sin^2{\theta})d\theta\\
    &= \frac{ab}{4}\int_{\theta= 0}^{2\pi}\left(a^2\cdot \frac{1+ \cos{2\theta}}{2}+ 2b^2\cdot \frac{1- \cos{2\theta}}{2}\right)d\theta\\
    &= \frac{ab}{4}\left(\frac{a^2}{2}\cdot 2\pi+ \frac{2b^2}{2}\cdot 2\pi\right)\\
    &= \frac{\pi ab(a^2+ 2b^2)}{4}
\end{align*}
    \item 条件$a^2+ 2b^2= 1$のもとで$V= \frac{\pi ab}{4}(a^2+ 2b^2)= \frac{\pi ab}{4}$の最大値を求める。\\
    相加相乗平均の関係を利用する。$a, b> 0$より$a^2, 2b^2> 0$。
\begin{equation*}
    1= a^2+ 2b^2\ge 2\sqrt{a^2\cdot 2b-2}= 2\sqrt{2}ab
\end{equation*}
    これより$ab\le \frac{1}{2\sqrt{2}}= \frac{\sqrt{2}}{4}$。したがって最大値$V_\mathrm{max}= \frac{\pi}{4}\cdot \frac{\sqrt{2}}{4}= \frac{\sqrt{2}}{16}\pi$
    \item 曲面$z= f(x, y)$の面積$A= \iint_D \sqrt{1+ \left(\frac{\partial z}{\partial x}\right)^2+ \left(\frac{\partial z}{\partial y}\right)^2}dxdy$で求められる。ここで$z= f(x, y)= x^2+ 2y^2$なので$\frac{\partial z}{\partial x}= 2x, \frac{\partial z}{\partial y}= 4y$。領域$x^2+ 2y^2\le 1$に合わせ、$x= r\cos{\theta}, y= \frac{1}{2}r\sin{\theta}$と置換する。ヤコビアンは$J= \frac{1}{2}r$。
\begin{align*}
    A
    &= \int_{\theta= 0}^{2\pi}\int_{r= 0}^1 \sqrt{1+ 4(x^2+ 2y^2)}\cdot \frac{1}{2}rdrd\theta\\
    &= \pi\left[\frac{1}{8}\cdot \frac{2}{3}(1+ 4r^2)^{3/2}\right]_{r= 0}^1\\
    &= \frac{\pi}{12}(5\sqrt{5}- 1)
\end{align*}
    \item 点$(x_0, y_0, z_0)$は$S$上にあるので$z= x_0^2+ 2y_0^2$を満たす関数を$F(x, y, z)= x^2+ 2y^2- z^2= 0$とおくと勾配ベクトルは
\begin{equation*}
    \nabla{F}= (2x, 4y, -1)
\end{equation*}
    点$(x_0, y_0, z_0)$における法線ベクトルは$(2x_0, 4y_0, -1)$なので、接平面の方程式は
\begin{equation*}
    2x_0(x- x_0)+ 4y(y- y_0)- (z- z_0)= 0
\end{equation*}
    これを変形すると接平面の方程式：$z- z_0= 2x_0(x- x_0)+ 4y_0(y- y_0)$または$z= 2x_0x+ 4y_0y- z$が得られる。
\end{enumerate}
\end{souanswer}
